\begin{table}[htbp]
\centering
\small
\caption{Confirmatory tests of H1–H3 with Holm–Bonferroni adjustment across the six-test family. All tests are paired within scenario.}
\label{tab:gmais-4}
\begin{tabular}{llrrrrrl}
\toprule
Hypothesis / contrast & Test & n & Statistic & p & p (Holm) & Effect size & Estimate [95\% CI] \\
\midrule
H1 | G=0: Baseline -> V-only & McNemar (exact) & 135 & 79.0 & < .0001 & < .0001 & odds ratio (discordant) = 53.000 & --- \\
H1 | G=1: G-only -> Full & McNemar (exact) & 135 & 60.0 & < .0001 & < .0001 & odds ratio (discordant) = 40.333 & --- \\
H2 | V=0: latency, Baseline -> G-only & Wilcoxon signed-rank & 135 & 3,700.0 & .0506 & .1519 & rank-biserial r = 0.194 & +103.876 [-0.362, +217.880] \\
H2 | V=1: latency, V-only -> Full & Wilcoxon signed-rank & 135 & 3,992.0 & .1891 & .3781 & rank-biserial r = 0.130 & +215.357 [-127.120, +505.132] \\
H2: token consumption, ungoverned -> governed & Wilcoxon signed-rank & 135 & 45.5 & < .0001 & < .0001 & rank-biserial r = 0.990 & +106.250 [+97.000, +116.500] \\
H3: V x G interaction on latency (ms) & Wilcoxon signed-rank & 135 & 4,402.0 & .6797 & .6797 & rank-biserial r = 0.041 & +79.944 [-312.409, +423.887] \\
\bottomrule
\end{tabular}
\par\vspace{2pt}\footnotesize\emph{Estimate is the Hodges–Lehmann shift with the exact distribution-free signed-rank interval read off the Walsh averages; for McNemar the effect size is the discordant-pair odds ratio (Haldane–Anscombe corrected). A degenerate interval indicates a differential that is constant across scenarios by construction --- see the note on the deployment cost model.}
\end{table}
